% Fixture for #138 — a control sheet shaped like a real one.
%
% This is what the generator will emit from the question bank: a pool of ten
% questions for a section range, of which each copy draws four at random, with
% the alternatives shuffled too (docs/design/2026-08-controles.md §C6).
%
% Content is Spanish because a control sheet is course material and a student
% reads it (apps/amc-worker/CLAUDE.md §Language). Macro names are AMC's English
% aliases — the .sty provides both those and the original French ones.
%
% The RUT grid is eight digits and carries no verifier digit: `\AMCcode` is
% digits-only and a DV can be K, and with an enrolled-student list the body
% alone identifies the student (§C5).

\documentclass[a4paper,11pt]{article}

\usepackage[utf8]{inputenc}
\usepackage[T1]{fontenc}
\usepackage[spanish]{babel}
\usepackage{multicol}
% lang=ES is not cosmetic: without it AMC labels every question "Question 1" in
% English, on a sheet a Spanish-speaking student reads. Measured in #138 S2 —
% the first compile shipped English labels past a green run, because nothing
% about a wrong-language label makes LaTeX fail.
\usepackage[box,completemulti,lang=ES]{automultiplechoice}

% Fixed seed so a test run is reproducible: the same seed must produce the same
% four questions in the same order for copy N, every time. Without this, a
% failure could never be told apart from a different draw.
\AMCrandomseed{1237}

\begin{document}

% ---------------------------------------------------------------------------
% The pool. In production these come from the published preguntas.json for the
% selected section range; here they stand in for it.
% ---------------------------------------------------------------------------

\element{clase}{
  \begin{question}{iteraciones}
    ¿Cuántas iteraciones hace, a lo más, la búsqueda binaria sobre un arreglo
    ordenado de $n$ elementos?
    \begin{choices}
      \correctchoice{Alrededor de $\log_2 n$}
      \wrongchoice{Alrededor de $n$}
      \wrongchoice{Alrededor de $n^2$}
      \wrongchoice{Siempre 2}
    \end{choices}
  \end{question}
}

\element{clase}{
  \begin{question}{requisito}
    ¿Qué condición debe cumplir el arreglo para poder aplicar búsqueda binaria?
    \begin{choices}
      \correctchoice{Estar ordenado}
      \wrongchoice{No tener elementos repetidos}
      \wrongchoice{Tener un número par de elementos}
      \wrongchoice{Estar lleno}
    \end{choices}
  \end{question}
}

\element{clase}{
  \begin{question}{lineal}
    En el peor caso, ¿cuántas comparaciones hace la búsqueda lineal sobre
    $n$ elementos?
    \begin{choices}
      \correctchoice{$n$}
      \wrongchoice{$\log_2 n$}
      \wrongchoice{$n/2$}
      \wrongchoice{1}
    \end{choices}
  \end{question}
}

\element{clase}{
  \begin{question}{tipo-primitivo}
    En Java, ¿cuál de estos es un tipo primitivo?
    \begin{choices}
      \correctchoice{\texttt{int}}
      \wrongchoice{\texttt{String}}
      \wrongchoice{\texttt{Integer}}
      \wrongchoice{\texttt{Object}}
    \end{choices}
  \end{question}
}

\element{clase}{
  \begin{question}{igualdad}
    En Java, ¿qué compara el operador \texttt{==} aplicado a dos objetos?
    \begin{choices}
      \correctchoice{Si son el mismo objeto}
      \wrongchoice{Si tienen el mismo contenido}
      \wrongchoice{Si son de la misma clase}
      \wrongchoice{Si ambos son distintos de \texttt{null}}
    \end{choices}
  \end{question}
}

\element{clase}{
  \begin{question}{arreglo-largo}
    En Java, ¿cómo se obtiene la cantidad de elementos de un arreglo
    \texttt{a}?
    \begin{choices}
      \correctchoice{\texttt{a.length}}
      \wrongchoice{\texttt{a.length()}}
      \wrongchoice{\texttt{a.size()}}
      \wrongchoice{\texttt{len(a)}}
    \end{choices}
  \end{question}
}

\element{clase}{
  \begin{question}{indice}
    ¿Cuál es el índice del primer elemento de un arreglo en Java?
    \begin{choices}
      \correctchoice{0}
      \wrongchoice{1}
      \wrongchoice{-1}
      \wrongchoice{Depende del tipo}
    \end{choices}
  \end{question}
}

\element{clase}{
  \begin{question}{referencia}
    Si \texttt{b = a} y ambos son referencias al mismo objeto, ¿qué ocurre al
    modificar un campo a través de \texttt{b}?
    \begin{choices}
      \correctchoice{El cambio se ve también a través de \texttt{a}}
      \wrongchoice{\texttt{a} conserva el valor anterior}
      \wrongchoice{Se lanza una excepción}
      \wrongchoice{Se crea una copia del objeto}
    \end{choices}
  \end{question}
}

\element{clase}{
  \begin{question}{descarte}
    En cada paso, la búsqueda binaria descarta:
    \begin{choices}
      \correctchoice{La mitad de los elementos que quedaban}
      \wrongchoice{Un elemento}
      \wrongchoice{Todos los mayores al buscado}
      \wrongchoice{Nada, solo avanza}
    \end{choices}
  \end{question}
}

\element{clase}{
  \begin{question}{peor-caso}
    Sobre 1000 elementos ordenados, ¿cuántas iteraciones hace la búsqueda
    binaria en el peor caso?
    \begin{choices}
      \correctchoice{Alrededor de 10}
      \wrongchoice{Alrededor de 100}
      \wrongchoice{Alrededor de 500}
      \wrongchoice{Alrededor de 1000}
    \end{choices}
  \end{question}
}

% ---------------------------------------------------------------------------
% The sheet. One per student, each drawing its own four questions.
% ---------------------------------------------------------------------------

\onecopy{5}{

  \noindent{\large\bf Control de entrada}\hfill
  \namefield{\fbox{\begin{minipage}{.45\linewidth}
        Nombre:

        \vspace*{4mm}\dotfill
        \vspace*{2mm}
      \end{minipage}}}

  \begin{center}
    \emph{Marca tu RUT sin el dígito verificador (8 dígitos).}

    \AMCcode{rut}{8}
  \end{center}

  \vspace{3mm}
  \noindent\emph{Marca una sola alternativa por pregunta.}
  \vspace{2mm}

  \shufflegroup{clase}
  \insertgroup[4]{clase}

  \AMCcleardoublepage
}

\end{document}
